% Part: computability
% Chapter: computability-theory
% Section: s-m-n

\documentclass[../../../include/open-logic-section]{subfiles}

\begin{document}

\olfileid{cmp}{thy}{smn}
\olsection{د $s$-$m$-$n$ قضيه}

\begin{explain}
راتلونکې قضيه د داسې دليل له مخې د ``$s$-$m$-$n$ قضيې'' په نوم
پېژندل کېږي چې لږ وروسته به څرګند شي۔ ستونزمنه برخه دا پوهېدل دي چې
قضيه په کره ډول څه وايي؛ کله چې ئې بيان درک کړئ، قضيه به بيخي څرګنده
وبرېښي۔
\end{explain}

\begin{thm}
\ollabel{thm:s-m-n}
  د طبيعي عددونو د هرې جوړې $n$ او~$m$ دپاره يوه بنسټيزه بازګشتي
  تابع~$s^m_n$ شته، داسې چې د هرې لړۍ
  $e$، $a_0$، \dots، $a_{m-1}$، $y_0$، \dots، $y_{n-1}$ دپاره لرو:
  \[
  \cfind{s^m_n(e, a_0, \dots, a_{m-1})}[n](y_0, \dots, y_{n-1}) \simeq
  \cfind{e}[m+n](a_0, \dots, a_{m-1}, y_0, \dots, y_{n-1}).
\]
\end{thm}

\begin{explain}
ګټوره ده چې $s^m_n$ پر \emph{پروګرامونو} د عمل کوونکې تابع په توګه
وګڼو۔ يعنې $s^m_n$ د يوې $(m+n)$-ځاييزې تابع پروګرام~$e$ او ثابت
ننوتونه $a_0$، \dots، $a_{m-1}$ اخلي؛ بيا د پاتې دليلونو د
$n$-ځاييزې تابع پروګرام $s^m_n(e, a_0, \dots, a_{m-1})$ راګرځوي۔
\iftag{TMs}{که~$e$ د يوه ټيورينګ ماشين تشريح وګڼئ، نو
$s^m_n(e, a_0, \dots, a_{m-1})$ هغه ټيورينګ ماشين دے چې د
$y_0$، \dots،~$y_{n-1}$ په ننوت کښې $a_0$، \dots،~$a_{m-1}$ د
ننوت د لړۍ په سر کښې ور زياتوي، او~$e$ چلوي۔ نو هره~$s^m_n$ يوازې
هغه بنسټيزه بازګشتي تابع ده چې د اړوند ټيورينګ ماشين کوډ پيدا کوي۔}{}
\textbf{د ژباړې د نسخې سمون (OLCMP-019):} سرچينې په همدې تشريح کښې
لومړے يو پروګرامي شاخص ټاکلے و، خو د راګرځېدونکي پروګرام په نښه کښې
ئې ناڅاپه بل متغير ليکلے و، او ورپسې د ټيورينګ ماشين تشريح هم د هماغه
بل متغير په نوم پيلوله۔ دلته دواړه ځايونه د قضيې له ټاکلي شاخص سره
يو شان ساتل شوي؛ د قضيې معادله نۀ ده بدله شوې۔
\end{explain}

\end{document}
